\begin{landscape}
\pagestyle{lscaped}
\begin{table}[!htbp]
\centering
\caption{Projected budget impact of Medicaid GLP-1 obesity coverage: covering states}
\label{tab:budget_impact_covering}
\begin{tabular}{@{} l l r rr rr @{}}
\toprule
 & & & \multicolumn{2}{c}{Average Post-Period} & \multicolumn{2}{c}{Event Time 9} \\
\cmidrule(lr){4-5} \cmidrule(lr){6-7}
State & Adoption & \shortstack{Enrollment\\(September 2025)} & \shortstack{Monthly\\Users} & \shortstack{Annual Net\\Spending (\$)} & \shortstack{Monthly\\Users} & \shortstack{Annual Net\\Spending (\$)} \\
\midrule
  California & 2023-Q1 & 11,978,256 & 42,238 & 477,682,515 & 98,341 & 1,080,410,097 \\
  Pennsylvania & 2023-Q1 & 2,757,077 & 9,722 & 109,949,852 & 22,636 & 248,681,764 \\
  North Carolina & 2024-Q3 & 2,535,360 & 8,940 & 101,107,969 & 20,815 & 228,683,420 \\
  Michigan & 2022-Q1 & 2,174,909 & 7,669 & 86,733,495 & 17,856 & 196,171,600 \\
  Virginia & 2023-Q1 & 1,512,223 & 5,332 & 60,306,149 & 12,415 & 136,398,905 \\
  Massachusetts & 2024-Q1 & 1,411,397 & 4,977 & 56,285,295 & 11,588 & 127,304,640 \\
  Tennessee & 2025-Q3 & 1,251,458 & 4,413 & 49,907,065 & 10,274 & 112,878,524 \\
  Minnesota & 2023-Q2 & 1,185,328 & 4,180 & 47,269,858 & 9,732 & 106,913,756 \\
  Missouri & 2024-Q4 & 1,133,687 & 3,998 & 45,210,459 & 9,308 & 102,255,861 \\
  Wisconsin & 2023-Q2 & 1,033,372 & 3,644 & 41,209,984 & 8,484 & 93,207,688 \\
  South Carolina & 2024-Q4 & 890,091 & 3,139 & 35,496,061 & 7,308 & 80,284,083 \\
  Mississippi & 2023-Q4 & 512,797 & 1,808 & 20,449,902 & 4,210 & 46,253,065 \\
  Kansas & 2024-Q2 & 333,192 & 1,175 & 13,287,410 & 2,736 & 30,053,123 \\
  Utah & 2025-Q3 & 302,064 & 1,065 & 12,046,052 & 2,480 & 27,245,452 \\
  Rhode Island & 2022-Q4 & 271,469 & 957 & 10,825,950 & 2,229 & 24,485,856 \\
  Delaware & 2023-Q1 & 230,166 & 812 & 9,178,822 & 1,890 & 20,760,424 \\
  New Hampshire & 2022-Q3 & 162,604 & 573 & 6,484,507 & 1,335 & 14,666,493 \\
\midrule
  \textbf{Subtotal (17 states)} &  & 29,675,450 & 104,643 & 1,183,431,344 & 243,635 & 2,676,654,748 \\
\bottomrule
\end{tabular}
\begin{minipage}{0.95\linewidth}
\vspace{0.5em}
\footnotesize
\textit{Notes.}
Projected monthly GLP-1 obesity drug users and annual Medicaid spending by state,
extrapolating per-enrollee ATT estimates from the 7-state enrollee-month weighted stacked
difference-in-differences design.
``Average post-period'' uses the mean ATT across event times 0--9 (approximately 2.5 years post-adoption).
``Event time 9'' uses the ATT at 9 quarters post-adoption, reflecting a more mature utilization level.
Users is the implied number of enrollees filling a GLP-1 obesity prescription in a given month.
Enrollment is total Medicaid enrollment in September 2025 from CMS monthly enrollment reports.
Net spending subtracts estimated statutory Medicaid Drug Rebate Program (MDRP) rebates from gross reimbursement (see Appendix~B).
Projections assume uniform per-enrollee effects across states and do not account for differences in
demographics, formulary design, or supplemental rebates.
\end{minipage}
\end{table}
\end{landscape}
